\problema{Find Median from Data Stream}{295}{src/04-heaps-topk/295-find-median-data-stream.cpp}{H}

Hay que diseñar una estructura \texttt{MedianFinder} que reciba números uno a
uno mediante \texttt{addNum} y en cualquier momento pueda devolver, con
\texttt{findMedian}, la mediana de todos los números recibidos hasta ese
instante. Si la cantidad de números vistos es impar, la mediana es el de en
medio; si es par, es el promedio de los dos de en medio.

\paragraph{La idea, con un ejemplo de la vida real.} Imagina una fila de
personas ordenadas de menor a mayor estatura, y quieres saber en todo
momento quién está justo en medio, aunque la fila crezca constantemente y
nunca la reordenes por completo. Una forma práctica es partir la fila en dos
mitades desde el principio: los más bajitos a la izquierda, los más altos a
la derecha, cuidando que ambos lados tengan casi el mismo número de
personas. La mediana siempre está en la frontera entre ambas mitades: la
persona más alta del lado izquierdo, la más baja del lado derecho, o el
promedio de ambas si las dos mitades tienen el mismo tamaño.

Para mantener esa frontera accesible en todo momento sin reordenar la fila
completa, guardamos la mitad izquierda en un \emph{max-heap} (para consultar
su máximo, es decir, el más alto del lado bajo, en tiempo constante) y la
mitad derecha en un \emph{min-heap} (para consultar su mínimo, el más bajo
del lado alto, también en tiempo constante).

\begin{center}
\ttfamily
\begin{tabular}{l}
addNum(1) -> stream: [1]\\
addNum(2) -> stream: [1,2] -> mediana = 1.5\\
addNum(3) -> stream: [1,2,3] -> mediana = 2\\
\end{tabular}
\end{center}

\paragraph{Cómo lo resuelve el código, paso a paso.}
\begin{enumerate}
  \item Mantenemos dos heaps: \texttt{izquierda} (max-heap, la mitad menor
        de los números vistos) y \texttt{derecha} (min-heap, la mitad
        mayor).
  \item Al insertar un número nuevo, decidimos a qué heap va comparándolo
        con la cima de \texttt{izquierda}: si es menor o igual (o si
        \texttt{izquierda} está vacía), va a \texttt{izquierda}; si no, va
        a \texttt{derecha}.
  \item Después de cada inserción, rebalanceamos: si algún heap le lleva más
        de un elemento de ventaja al otro, movemos su cima al heap
        contrario. Esto mantiene la diferencia de tamaños en, a lo más, 1.
  \item Para consultar la mediana: si ambos heaps tienen el mismo tamaño, es
        el promedio de sus dos cimas; si no, es la cima del heap más grande.
\end{enumerate}
El invariante clave que se mantiene todo el tiempo es que \emph{todo}
elemento de \texttt{izquierda} es menor o igual que \emph{todo} elemento de
\texttt{derecha}. Como cada inserción respeta ese orden y el rebalanceo solo
mueve el elemento frontera de un heap al otro, ese invariante nunca se
rompe, y por eso la mediana siempre puede leerse directamente de las cimas.

\lstinputlisting{src/04-heaps-topk/295-find-median-data-stream.cpp}

\complejidad{$O(\log n)$ por \texttt{addNum}, $O(1)$ por \texttt{findMedian}}{$O(n)$}

\paragraph{¿Qué significa esa complejidad?} Cada número que entra puede
costar hasta dos operaciones de heap (una inserción y, en el peor caso, un
rebalanceo), cada una $O(\log n)$ porque los heaps van creciendo con el
stream. En cambio, consultar la mediana no requiere ningún trabajo extra:
las cimas de ambos heaps ya están calculadas, así que es simplemente leerlas.

\paragraph{Consejos para la entrevista (Amazon).} Este es un problema
clásico de diseño de estructura de datos con heaps, y suelen evaluar si
reconoces cuándo un problema deja de ser ``calcula una respuesta'' y pasa a
ser ``mantén una respuesta actualizable eficientemente'' conforme llegan
datos nuevos. El error más común es rebalancear mal: olvidar el caso en que
\texttt{izquierda} está vacía al decidir dónde insertar el primer número, o
permitir que la diferencia de tamaños crezca más de 1 (rompiendo la
garantía de $O(1)$ en \texttt{findMedian}). Vale la pena aclarar en voz alta
por qué se usa \texttt{<=} y no \texttt{<} al decidir el heap destino: con
duplicados, cualquiera de las dos convenciones funciona mientras se use de
forma consistente.
